\documentclass[11pt]{article}
\usepackage[margin=0.85in]{geometry}
\usepackage{amsmath,amssymb}
\title{Torus configuration rings generated in degree one}
\author{The Last Math Competition, conjecture 00000004098}
\date{}
\begin{document}
\maketitle

\paragraph{One particle, integral coefficients.}
For one particle the ordered and unordered configuration spaces are
both \(T^2=S^1\times S^1\). The integral Künneth theorem gives
\[
  H^*(T^2;\mathbb Z)=\Lambda_{\mathbb Z}(a,b), \qquad |a|=|b|=1.
\]
The classes \(a,b\) generate, with \(ab\) giving the degree-two
class. Any homogeneous generating set must span the rank-two
degree-one group in degree one. Thus the minimum number of
generators is \(2\), whereas \(n+2=3\) for \(n=1\). The least
possible maximal generating degree is \(1\). This already refutes
both numerical assertions.

\paragraph{Two ordered particles, integral coefficients.}
Writing the torus additively, the map
\[
  F_2(T^2)\longrightarrow T^2\times(T^2\setminus\{0\}),\qquad
  (x,y)\longmapsto (x,y-x)
\]
is a homeomorphism with inverse \((x,u)\mapsto (x,x+u)\).
The punctured torus retracts to a wedge of two circles. For example,
take a square fundamental domain with the puncture at its center
and retract the punctured square radially to its boundary, fixing that
boundary before its edge identifications.
Consequently the integral Künneth isomorphism for the product gives
\[
  H^*(F_2(T^2);\mathbb Z)=
  \Lambda_{\mathbb Z}(a,b,c,d)/(cd),
  \qquad |a|=|b|=|c|=|d|=1.
\]
Here \(a,b\) are from the torus factor and \(c,d\) are the
restrictions of the two circle classes to the punctured torus factor.
In particular, degree-one generators suffice even for two particles.

\paragraph{Two unordered particles, rational coefficients.}
The deck involution on the ordered space in the above coordinates is
\((x,u)\mapsto(x+u,-u)\). Addition on each torus circle pulls the
degree-one class back to the sum from the two factors, and inversion
negates that class. Thus the deck pullback is
\[
  a\longmapsto a+c,\qquad b\longmapsto b+d,\qquad
  c\longmapsto-c,\qquad d\longmapsto-d .
\]
By transfer, the rational cohomology of \(F_2(T^2)/S_2\) identifies
with the invariant subring of the ordered cohomology
\cite[Proposition 3G.1]{hatcher}. Set \(A=2a+c\) and \(B=2b+d\).
The four classes \(A,B,c,d\) are a rational basis in degree one, with
\(A,B\) invariant and \(c,d\) negated. In this basis the sole exterior
relation beyond anticommutation and squares is still \(cd=0\).
Every surviving monomial involving \(c\) or \(d\) contains exactly
one of them and is negated. Hence the invariant subring is
\(\Lambda_{\mathbb Q}(A,B)\), generated by two classes of degree one.

\paragraph{Lean verification.}
The Lean project checks the exterior monomial products and generating
degrees for the one- and two-particle calculations. It verifies the
deck involution and the invariant monomial basis over characteristic
zero. The identifications with configuration-space cohomology follow
from the topology, the Künneth theorem and transfer as above.

\begin{thebibliography}{1}
\bibitem{hatcher}
A. Hatcher, \emph{Algebraic Topology}, Cambridge University Press,
2002, Section 3.G.
\end{thebibliography}
\end{document}
